% ===========================================================================
%  Chapitre 10 — Généralités sur les suites numériques
%  RAPE, quatrième période — hors tableau des contenus du PE 2026
% ===========================================================================
\bmtchapitre{Généralités sur les suites numériques}
  {Reconnaître une suite donnée sous forme explicite ou récurrente, calculer
   ses premiers termes, étudier son sens de variation, et déterminer sa
   limite : convergence ou divergence.}
  {Suites numériques \textbullet\ environ 6 heures}

Une suite est une liste infinie de nombres, numérotée par les entiers. Dit
comme cela, l'objet paraît maigre. Il ne l'est pas : une suite décrit tout ce
qui se répète pas à pas — un capital qui produit des intérêts chaque année,
une population qui augmente d'un pourcentage fixe, un effectif scolaire relevé
session après session.

Ce chapitre installe le vocabulaire et les deux gestes de base : calculer des
termes, et dire où la suite va quand $n$ devient très grand. Le chapitre
suivant traitera des deux familles qui font l'exercice 1 du baccalauréat.

% ---------------------------------------------------------------------------
\begin{revision}
Les chapitres 4, 7 et 8 suffisent.

\begin{enumerate}[label=\textbf{\arabic*.}, leftmargin=*, itemsep=0.35em]
  \item Calculer $f(3)$ pour $f(x) = \dfrac{x-1}{x}$.
  \item Calculer $\limpinf \dfrac{2x+1}{x}$.
  \item Que vaut $e^{0}$ ? $\ln 1$ ?
  \item Comparer $\dfrac{1}{100}$ et $\dfrac{1}{1\,000}$.
  \item Que devient $\left(\frac12\right)^{n}$ quand $n$ devient très
        grand ?
  \item Résoudre $2x-6 \geq 0$.
\end{enumerate}
\end{revision}

\begin{automatismes}
Sans calculatrice.

\begin{enumerate}[label=\textbf{\alph*)}, leftmargin=*, itemsep=0.25em]
  \item $u_{n} = 2n+1$ : donner $u_{0}$ et $u_{3}$
  \item $u_{n} = n^{2}-n$ : donner $u_{4}$
  \item $u_{0} = 5$ et $u_{n+1} = u_{n}-2$ : donner $u_{1}$ et $u_{2}$
  \item $u_{0} = 1$ et $u_{n+1} = 3u_{n}$ : donner $u_{2}$
  \item $\limn \dfrac1n$
  \item $\limn \left(\frac13\right)^{n}$
  \item $\limn 2^{n}$
  \item $u_{n} = (-1)^{n}$ : donner $u_{10}$ et $u_{11}$
\end{enumerate}
\end{automatismes}

\begin{corrige}[automatismes]
a) $1$ et $7$ \quad b) $12$ \quad c) $3$ et $1$ \quad d) $9$ \quad
e) $0$ \quad f) $0$ \quad g) $+\infty$ \quad h) $1$ et $-1$.
\end{corrige}

% ===========================================================================
\section{Définition et modes de génération}

\begin{definition}[suite numérique]
Une \textbf{suite numérique} est une fonction définie sur $\N$ (ou sur une
partie de $\N$) et à valeurs réelles. L'image de l'entier $n$ se note
$u_{n}$ et s'appelle le \textbf{terme de rang} $n$. La suite elle-même se
note $\left(u_{n}\right)$.
\end{definition}

\begin{avertissement}
$u_{n}$ est un \emph{nombre} — le terme de rang $n$ ; $\left(u_{n}\right)$
est la \emph{suite} tout entière. Les parenthèses ne sont pas décoratives :
« la suite $u_{n}$ est croissante » n'a pas de sens, seule
« la suite $\left(u_{n}\right)$ est croissante » en a un.
\end{avertissement}

\begin{propriete}[deux façons de définir une suite]
\begin{itemize}[leftmargin=1.3em, itemsep=0.2em]
  \item \textbf{Forme explicite} : $u_{n} = f(n)$. Chaque terme se calcule
        directement à partir de son rang. Par exemple $u_{n} = 3n-1$ donne
        $u_{100} = 299$ sans calculer les précédents.
  \item \textbf{Forme récurrente} : un premier terme et une relation
        $u_{n+1} = g\!\left(u_{n}\right)$. Chaque terme se déduit du
        précédent, et l'on ne peut pas sauter d'étape.
\end{itemize}
\end{propriete}

\begin{exemple}[la même question, deux formes]
\textbf{a)} $u_{n} = \dfrac{n-1}{n}$ pour $n \geq 1$. Alors
$u_{1} = 0$, $u_{2} = \frac12$, $u_{3} = \frac23$. On peut aussi écrire
directement $u_{3n+1} = \dfrac{3n+1-1}{3n+1} = \dfrac{3n}{3n+1}$.

\textbf{b)} $u_{0} = 4$ et $u_{n+1} = \frac34u_{n}-1$. Alors
$u_{1} = 3-1 = 2$ et $u_{2} = \frac32-1 = \frac12$. Pour obtenir
$u_{100}$, il faudrait cent calculs : c'est tout l'intérêt des méthodes du
chapitre suivant, qui donnent une forme explicite.
\end{exemple}

\begin{visualisation}[les termes d'une suite sur un axe]
\begin{center}
\begin{tikzpicture}
\begin{axis}[
  width = 11.4cm, height = 5.2cm,
  axis lines = left, xlabel = {$n$}, ylabel = {$u_{n}$},
  xmin = -0.4, xmax = 10.5, ymin = -0.2, ymax = 4.6,
  xtick = {0,1,2,3,4,5,6,7,8,9,10}, ytick = {1,2,3,4},
  axis line style = {bmtGris},
  tick label style = {font=\scriptsize, color=bmtGris},
  grid = both, grid style = {bmtGrisClair, line width=0.3pt}, clip = true,
]
  \addplot[only marks, mark=*, mark size=1.8pt, bmtPalissandre,
           domain=0:10, samples=11] {2+2*(0.72)^x};
  \addplot[bmtLaterite, line width=0.9pt, dash pattern=on 3pt off 2pt,
           domain=-0.4:10.5] {2};
  \node[bmtLaterite, font=\scriptsize, anchor=south east]
    at (axis cs:10.3,2.05) {$\ell = 2$};
\end{axis}
\end{tikzpicture}
\end{center}

Une suite ne se trace pas comme une courbe : elle se marque par des points
isolés, un par rang. Ici les termes descendent et s'accumulent au-dessus de
la valeur $2$ sans jamais l'atteindre : la suite \emph{converge} vers $2$.
\end{visualisation}

\begin{entrainement}[premiers termes]
\exo[\niveau{1}] $u_{n} = 4n-3$ : calculer $u_{0}$, $u_{1}$ et $u_{10}$.

\exo[\niveau{1}] $u_{0} = 2$ et $u_{n+1} = 2u_{n}+1$ : calculer $u_{1}$,
$u_{2}$ et $u_{3}$.

\exo[\niveau{2}] $u_{n} = \dfrac{n+2}{n+1}$ : calculer $u_{0}$ et $u_{5}$,
puis exprimer $u_{2n}$ en fonction de $n$.

\exo[\niveau{2}] $u_{0} = 1$ et $u_{n+1} = \dfrac{u_{n}+3}{4}$ : calculer
$u_{1}$ et $u_{2}$.

\exo[\niveau{3}] $u_{n} = e^{2-n}$ : calculer $u_{0}$, $u_{1}$ et $u_{2}$ en
fonction de $e$, puis exprimer $\dfrac{u_{n+1}}{u_{n}}$.
\end{entrainement}

\begin{corrige}[premiers termes]
\textbf{1.} $-3$ ; $1$ ; $37$.
\textbf{2.} $5$ ; $11$ ; $23$.
\textbf{3.} $2$ ; $\frac76$ ; $u_{2n} = \frac{2n+2}{2n+1}$.
\textbf{4.} $u_{1} = 1$ et $u_{2} = 1$ — la suite est constante, car $1$ est
le point fixe de la relation.
\textbf{5.} $e^{2}$ ; $e$ ; $1$ ; et
$\frac{u_{n+1}}{u_{n}} = \frac{e^{1-n}}{e^{2-n}} = e^{-1}$.
\end{corrige}

% ===========================================================================
\section{Sens de variation}

\begin{definition}[suite croissante, décroissante]
La suite $\left(u_{n}\right)$ est \textbf{croissante} si
$u_{n+1} \geq u_{n}$ pour tout $n$ ; \textbf{décroissante} si
$u_{n+1} \leq u_{n}$ pour tout $n$ ; \textbf{constante} si
$u_{n+1} = u_{n}$ pour tout $n$.
\end{definition}

\begin{methode}{étudier le sens de variation}
Trois techniques, à choisir selon l'écriture de la suite.
\begin{enumerate}[label=\textbf{\arabic*.}, leftmargin=*, itemsep=0.15em]
  \item \emph{La différence.} Calculer $u_{n+1}-u_{n}$ et étudier son signe.
        C'est la méthode par défaut, et la seule qui marche toujours.
  \item \emph{Le quotient.} Si tous les termes sont strictement positifs,
        comparer $\dfrac{u_{n+1}}{u_{n}}$ à $1$. Commode pour les suites
        écrites avec des puissances ou une exponentielle.
  \item \emph{La fonction associée.} Si $u_{n} = f(n)$ et si $f$ est
        croissante sur $\intervallefo{0}{+\infty}$, la suite est croissante.
        Attention : cette méthode ne s'applique \textbf{pas} aux suites
        définies par récurrence.
\end{enumerate}
\end{methode}

\begin{exemple}[trois suites, trois techniques]
\textbf{a)} $u_{n} = n^{2}-3n$. Alors
\[
  u_{n+1}-u_{n} = \left[(n+1)^{2}-3(n+1)\right]-\left[n^{2}-3n\right]
  = 2n-2 ,
\]
négatif pour $n = 0$, nul pour $n = 1$, positif ensuite : la suite décroît
puis croît.

\textbf{b)} $u_{n} = \dfrac{3^{n}}{2}$, à termes strictement positifs.
Alors $\dfrac{u_{n+1}}{u_{n}} = 3 > 1$ : la suite est croissante.

\textbf{c)} $u_{n} = \dfrac{n-1}{n}$ pour $n \geq 1$. La fonction
$f(x) = \frac{x-1}{x} = 1-\frac1x$ a pour dérivée $\frac{1}{x^{2}}>0$ : elle
est croissante, donc la suite aussi.
\end{exemple}

\begin{entrainement}[sens de variation]
\exo[\niveau{1}] $u_{n} = 5-2n$ : étudier le sens de variation.

\exo[\niveau{2}] $u_{n} = n^{2}+n$ : montrer que la suite est croissante.

\exo[\niveau{2}] $u_{n} = \left(\frac25\right)^{n}$ : étudier le sens de
variation par le quotient.

\exo[\niveau{2}] $u_{n} = \dfrac{1}{n+1}$ : étudier le sens de variation.

\exo[\niveau{3}] $u_{n} = e^{3n-1}$ : montrer que la suite est croissante.

\exo[\niveau{3}] $u_{0} = 0$ et $u_{n+1} = \sqrt{u_{n}+2}$ : calculer les
trois premiers termes et conjecturer le sens de variation.
\end{entrainement}

\begin{corrige}[sens de variation]
\textbf{1.} $u_{n+1}-u_{n} = -2<0$ : décroissante.
\textbf{2.} $u_{n+1}-u_{n} = 2n+2>0$ : croissante.
\textbf{3.} $\frac{u_{n+1}}{u_{n}} = \frac25<1$ et les termes sont positifs :
décroissante.
\textbf{4.} $u_{n+1}-u_{n} = \frac{1}{n+2}-\frac{1}{n+1}
= \frac{-1}{(n+1)(n+2)}<0$ : décroissante.
\textbf{5.} $\frac{u_{n+1}}{u_{n}} = e^{3}>1$, termes positifs :
croissante.
\textbf{6.} $u_{0} = 0$, $u_{1} = \sqrt2 \approx 1{,}41$,
$u_{2} = \sqrt{2+\sqrt2} \approx 1{,}85$ : la suite semble croissante.
\end{corrige}

% ===========================================================================
\section{Limite d'une suite}

\begin{definition}[convergence, divergence]
La suite $\left(u_{n}\right)$ \textbf{converge} vers le réel $\ell$ si
$u_{n}$ devient aussi proche que l'on veut de $\ell$ dès que $n$ est assez
grand ; on note $\limn u_{n} = \ell$. Une suite qui ne converge pas est dite
\textbf{divergente} : soit parce qu'elle tend vers $\pm\infty$, soit parce
qu'elle n'a pas de limite du tout.
\end{definition}

\begin{propriete}[limites de référence]
\[
  \limn \frac1n = 0, \qquad \limn \frac{1}{n^{2}} = 0, \qquad
  \limn n = +\infty, \qquad \limn n^{2} = +\infty .
\]
Et, pour un réel $q$ :
\[
  \limn q^{n} =
  \begin{cases}
    0 & \text{si } -1<q<1,\\
    1 & \text{si } q = 1,\\
    +\infty & \text{si } q>1,\\
    \text{pas de limite} & \text{si } q \leq -1 .
  \end{cases}
\]
\end{propriete}

\begin{remarque}
Cette dernière propriété est la plus utilisée de tout le chapitre : c'est
elle qui donne la limite de toute suite géométrique, et donc la réponse de
la dernière question de l'exercice 1 à presque chaque session.
\end{remarque}

\begin{propriete}[suites écrites avec $e$ ou $\ln$]
\[
  \limn e^{-n} = 0, \qquad \limn e^{n} = +\infty, \qquad
  \limn \ln n = +\infty .
\]
Plus généralement, pour $a$ réel non nul, $e^{an}$ se comporte comme
$\left(e^{a}\right)^{n}$ : c'est une suite géométrique de raison $e^{a}$, et
la règle sur $q^{n}$ s'applique, avec $q = e^{a}$.
\end{propriete}

\begin{exemple}[trois limites]
\textbf{a)} $u_{n} = 3+\dfrac{2}{n}$ : la limite vaut $3$, la suite
converge.

\textbf{b)} $u_{n} = 4\left(\frac13\right)^{n}$ : comme
$\left|\frac13\right|<1$, la limite vaut $0$.

\textbf{c)} $u_{n} = e^{2-n} = e^{2}\times\left(e^{-1}\right)^{n}$ : la
raison $e^{-1} \approx 0{,}37$ est comprise entre $-1$ et $1$, donc la
limite vaut $0$.
\end{exemple}

\begin{erreurclassique}
Une suite dont les termes deviennent de plus en plus grands n'est pas
« convergente vers l'infini » : elle est \textbf{divergente}. Le mot
« converge » est réservé aux limites \emph{finies}. Cette distinction est
exigée dans les corrigés officiels.
\end{erreurclassique}

\begin{entrainement}[limites]
\exo[\niveau{1}] $u_{n} = 2-\dfrac3n$

\exo[\niveau{1}] $u_{n} = n^{2}+1$

\exo[\niveau{2}] $u_{n} = \left(\frac34\right)^{n}$

\exo[\niveau{2}] $u_{n} = 5\left(\frac54\right)^{n}$

\exo[\niveau{2}] $u_{n} = \dfrac{2n+1}{n+3}$

\exo[\niveau{3}] $u_{n} = e^{-2n}$

\exo[\niveau{3}] $u_{n} = 3+\left(-\frac12\right)^{n}$

\exo[\niveau{3}] $u_{n} = (-1)^{n}+1$
\end{entrainement}

\begin{corrige}[limites]
\textbf{1.} $2$ (convergente).
\textbf{2.} $+\infty$ (divergente).
\textbf{3.} $0$.
\textbf{4.} $+\infty$, la raison dépassant $1$.
\textbf{5.} $\frac{2n}{n} = 2$.
\textbf{6.} $\left(e^{-2}\right)^{n} \to 0$.
\textbf{7.} $3$, car $\left(-\frac12\right)^{n} \to 0$.
\textbf{8.} Pas de limite : la suite vaut alternativement $2$ et $0$.
\end{corrige}

\begin{qcm}
\exo Pour $u_{n} = 7-n$, la suite est :
\textbf{a)} croissante \quad \textbf{b)} décroissante \quad
\textbf{c)} constante \quad \textbf{d)} ni l'un ni l'autre \reponseqcm{b}

\exo $\limn \left(\frac74\right)^{n}$ vaut :
\textbf{a)} $0$ \quad \textbf{b)} $1$ \quad \textbf{c)} $+\infty$ \quad
\textbf{d)} $\frac74$ \reponseqcm{c}

\exo Une suite définie par $u_{n+1} = g(u_{n})$ est dite :
\textbf{a)} explicite \quad \textbf{b)} récurrente \quad
\textbf{c)} arithmétique \quad \textbf{d)} bornée \reponseqcm{b}

\exo Si $\limn u_{n} = +\infty$, la suite est :
\textbf{a)} convergente \quad \textbf{b)} divergente \quad
\textbf{c)} constante \quad \textbf{d)} décroissante \reponseqcm{b}

\exo $u_{n} = e^{1-n}$ : cette suite est :
\textbf{a)} croissante de limite $+\infty$ \quad
\textbf{b)} décroissante de limite $0$ \quad
\textbf{c)} constante \quad \textbf{d)} sans limite \reponseqcm{b}
\end{qcm}

% ===========================================================================
\begin{annales}
Chaque session ouvre par un exercice de suites. Les extraits ci-dessous en
sont les premières questions — celles qui ne demandent que le vocabulaire et
les calculs de ce chapitre. Les exercices complets figurent au chapitre
suivant.

\exoann{Baccalauréat, Madagascar, série A, session 2010 — exercice 2,
question 1}
$\left(U_{n}\right)_{n \in \N^{*}}$ est la suite numérique définie par
$U_{n} = \dfrac{n-1}{n}$.
\begin{enumerate}[label=\textbf{\alph*)}, leftmargin=*, itemsep=0.15em]
  \item Calculer les trois premiers termes de $\left(U_{n}\right)$.
  \item Exprimer $U_{3n+1}$ en fonction de $n$.
\end{enumerate}

\exoann{Baccalauréat, Madagascar, série A, session 2013 — exercice 1,
question 1}
Soit $\left(U_{n}\right)$ définie par $U_{0} = 4$ et
$U_{n+1} = \dfrac34 U_{n}-1$ pour tout $n \in \N$. Calculer $U_{1}$ et
$U_{2}$.

\exoann{Baccalauréat, Madagascar, série A, session 2014 — exercice 1,
question 1}
Soit $\left(U_{n}\right)$ définie par $U_{0} = -4$ et
$U_{n+1} = \dfrac13 U_{n}-4$. Calculer $U_{1}$ et $U_{2}$.

\exoann{Baccalauréat, Madagascar, série A, session 2026 — exercice 1,
question 1}
Soit $\left(U_{n}\right)$ définie par $U_{0} = 5$ et
$U_{n+1} = \dfrac{U_{n}+3}{4}$. Calculer $U_{1}$ et $U_{2}$.

\exoann{Baccalauréat, Madagascar, série A, session 2017 — exercice 1,
question 1.a)}
Pour tout entier naturel $n$, on pose $U_{n} = \left(\frac34\right)^{n}$ et
$V_{n} = \ln\left(U_{n}\right)$. Calculer $U_{0}$, $U_{1}$, $V_{0}$ et
$V_{1}$.

\exoann{Baccalauréat, Madagascar, série A, session 2020 — exercice 1,
question 1}
$\left(V_{n}\right)$ est la suite définie par
$V_{n} = 3\left(\frac58\right)^{n}$.
\begin{enumerate}[label=\textbf{\alph*)}, leftmargin=*, itemsep=0.15em]
  \item Calculer $V_{0}$ et $V_{1}$.
  \item Exprimer $V_{n+1}$ en fonction de $V_{n}$.
  \item Calculer $\limn V_{n}$.
\end{enumerate}
\end{annales}

\begin{corrige}[annales]
\textbf{1. a)} $U_{1} = 0$, $U_{2} = \frac12$, $U_{3} = \frac23$.
\textbf{b)} $U_{3n+1} = \frac{3n}{3n+1}$.

\textbf{2.} $U_{1} = 3-1 = 2$ et $U_{2} = \frac32-1 = \frac12$.

\textbf{3.} $U_{1} = -\frac43-4 = -\frac{16}{3}$ et
$U_{2} = -\frac{16}{9}-4 = -\frac{52}{9}$.

\textbf{4.} $U_{1} = \frac{5+3}{4} = 2$ et $U_{2} = \frac{2+3}{4}
= \frac54$.

\textbf{5.} $U_{0} = 1$, $U_{1} = \frac34$, $V_{0} = \ln1 = 0$ et
$V_{1} = \ln\frac34 = \ln3-\ln4 \approx -0{,}29$.

\textbf{6. a)} $V_{0} = 3$ et $V_{1} = \frac{15}{8}$.
\textbf{b)} $V_{n+1} = 3\left(\frac58\right)^{n+1}
= \frac58\times3\left(\frac58\right)^{n} = \frac58 V_{n}$.
\textbf{c)} Comme $\left|\frac58\right|<1$, $\limn V_{n} = 0$.
\end{corrige}

% ===========================================================================
\begin{vraifaux}
\exo Une suite croissante est nécessairement divergente.

\exo Si $u_{n} = f(n)$ avec $f$ décroissante sur
$\intervallefo{0}{+\infty}$, alors $\left(u_{n}\right)$ est décroissante.

\exo Une suite qui tend vers $+\infty$ est convergente.

\exo Deux suites peuvent avoir les mêmes premiers termes sans être égales.

\exo $\limn \left(-\frac34\right)^{n} = 0$.
\end{vraifaux}

\begin{corrige}[vrai ou faux]
\textbf{1.} Faux — $u_{n} = 1-\frac1n$ croît et converge vers $1$.
\textbf{2.} Vrai — c'est la méthode de la fonction associée.
\textbf{3.} Faux — elle est divergente ; « converger » suppose une limite
finie.
\textbf{4.} Vrai — connaître dix termes ne dit rien du onzième ; seule la
définition compte.
\textbf{5.} Vrai — la valeur absolue de la raison, $\frac34$, est
inférieure à $1$.
\end{corrige}

% ===========================================================================
\begin{situation}[la citerne du collège]
La citerne d'un collège contient $8\,000$ litres d'eau au premier jour du
trimestre. Chaque jour, les élèves en prélèvent $10\,\%$ du contenu, puis la
pompe en ajoute $500$ litres. On note $u_{n}$ le nombre de litres présents
au début du $n$-ième jour, avec $u_{0} = 8\,000$.

\begin{enumerate}[label=\textbf{\arabic*.}, leftmargin=*, itemsep=0.3em]
  \item Justifier que $u_{n+1} = 0{,}9\,u_{n}+500$.
  \item Calculer $u_{1}$, $u_{2}$ et $u_{3}$.
  \item La suite semble-t-elle croissante ou décroissante ?
  \item Montrer que si $u_{n} = 5\,000$, alors $u_{n+1} = 5\,000$.
        Comment appelle-t-on une telle valeur ?
  \item Que peut-on conjecturer sur la limite de la suite ?
\end{enumerate}
\end{situation}

\begin{corrige}[la citerne du collège]
\textbf{1.} Prélever $10\,\%$ laisse $90\,\%$, soit $0{,}9\,u_{n}$ ; on
ajoute ensuite $500$ litres.

\textbf{2.} $u_{1} = 7\,700$, $u_{2} = 7\,430$, $u_{3} = 7\,187$.

\textbf{3.} Décroissante : chaque jour, le prélèvement dépasse l'apport.

\textbf{4.} $0{,}9\times5\,000+500 = 4\,500+500 = 5\,000$. Une telle valeur,
qui se reproduit d'elle-même, s'appelle un \textbf{point fixe} de la
relation.

\textbf{5.} La suite décroît et reste au-dessus de $5\,000$ : on peut
conjecturer qu'elle converge vers $5\,000$. Le chapitre suivant permettra de
le démontrer.
\end{corrige}

% ===========================================================================
\begin{jemeteste}
Vingt minutes.

\begin{enumerate}[label=\textbf{\arabic*.}, leftmargin=*, itemsep=0.3em]
  \item $u_{n} = 3n-7$ : calculer $u_{0}$, $u_{4}$ et $u_{n+1}$.
  \item $u_{0} = 6$ et $u_{n+1} = \frac12u_{n}+1$ : calculer $u_{1}$ et
        $u_{2}$.
  \item Étudier le sens de variation de $u_{n} = \dfrac{2}{n+3}$.
  \item Calculer $\limn \left(\frac49\right)^{n}$ et
        $\limn \left(\frac94\right)^{n}$.
  \item $u_{n} = e^{2-3n}$ : montrer que la suite est décroissante et donner
        sa limite.
\end{enumerate}
\end{jemeteste}

\begin{corrige}[je me teste]
\textbf{1.} $-7$ ; $5$ ; $u_{n+1} = 3n-4$.
\textbf{2.} $u_{1} = 4$ et $u_{2} = 3$.
\textbf{3.} $u_{n+1}-u_{n} = \frac{2}{n+4}-\frac{2}{n+3}
= \frac{-2}{(n+3)(n+4)}<0$ : décroissante.
\textbf{4.} $0$ et $+\infty$.
\textbf{5.} $\frac{u_{n+1}}{u_{n}} = e^{-3}<1$ avec des termes positifs :
décroissante ; et $u_{n} = e^{2}\left(e^{-3}\right)^{n} \to 0$.
\end{corrige}

% ===========================================================================
\begin{commeaubac}
\bmtsource{Exercice original au format de l'épreuve — série A, options A1 et
A2 \textbullet\ 5 points}
On considère la suite $\left(U_{n}\right)$ définie par $U_{0} = 9$ et, pour
tout entier naturel $n$,
\[
  U_{n+1} = \frac{U_{n}+6}{2}.
\]

\begin{enumerate}[label=\textbf{\arabic*.}, leftmargin=*, itemsep=0.25em]
  \item Calculer $U_{1}$, $U_{2}$ et $U_{3}$. \bareme{1 pt}
  \item Montrer que si $U_{n} = 6$, alors $U_{n+1} = 6$. \bareme{0,5 pt}
  \item Calculer $U_{n+1}-U_{n}$ en fonction de $U_{n}$, puis montrer que la
        suite est décroissante tant que $U_{n}>6$. \bareme{1,5 pt}
  \item On pose $V_{n} = U_{n}-6$. Calculer $V_{0}$, $V_{1}$ et $V_{2}$.
        \bareme{1 pt}
  \item Conjecturer la limite de $\left(U_{n}\right)$ et justifier en une
        phrase. \bareme{1 pt}
\end{enumerate}
\end{commeaubac}

\begin{corrige}[comme au baccalauréat]
\textbf{1.} $U_{1} = \frac{15}{2} = 7{,}5$ ; $U_{2} = \frac{13{,}5}{2}
= 6{,}75$ ; $U_{3} = 6{,}375$.

\textbf{2.} $\frac{6+6}{2} = 6$ : la valeur $6$ est un point fixe.

\textbf{3.} $U_{n+1}-U_{n} = \frac{U_{n}+6}{2}-U_{n}
= \frac{6-U_{n}}{2}$. Cette différence est négative dès que $U_{n}>6$ :
la suite est alors décroissante.

\textbf{4.} $V_{0} = 3$, $V_{1} = 1{,}5$, $V_{2} = 0{,}75$ : chaque terme
est la moitié du précédent.

\textbf{5.} Les termes $V_{n}$ semblent être ceux d'une suite géométrique de
raison $\frac12$, donc de limite $0$ ; par conséquent $U_{n} = V_{n}+6$
semble converger vers $6$. C'est exactement la démarche que le chapitre
suivant rendra rigoureuse.
\end{corrige}

\begin{notepourlaclasse}
Six heures pour les deux chapitres de la partie, dont deux ici : ce chapitre
n'est qu'un vestibule, l'essentiel est au chapitre 11. Ne vous attardez ni
sur la définition formelle de la limite, ni sur les théorèmes de comparaison,
qui ne sont pas au programme de la série.

Le point à travailler est l'écriture $u_{n} = a\,q^{n}$ pour les suites
écrites avec une exponentielle : $e^{2-n} = e^{2}\times\left(e^{-1}\right)^{n}$.
Les sujets en posent une à peu près une session sur deux, et c'est la seule
difficulté technique de ce chapitre.

Le bloc « comme au baccalauréat » prépare directement l'exercice 1 du
chapitre suivant : il introduit la suite auxiliaire $V_{n} = U_{n}-6$ sans
la nommer, pour que le mécanisme soit déjà familier lorsqu'il sera énoncé.
\end{notepourlaclasse}
